% ---
% title: Hinge loss: An example
% date: 2026-03-30
% tag: ML
% excerpt: 一个带 margin 的排序损失：正样本要靠近正确类别，负样本至少隔开 m；顺带和 triplet loss、交叉熵做个对比。
% ---
\[
\mathcal{L}_{hinge} = \sum_{k=1}^{|\mathcal{P}|} \max\bigl(\text{dist}(A(\hat{e}_c), f_{text}(t_c)) - \text{dist}(A(\hat{e}_c), f_{text}(t_{\cancel{c}})) + m,\ 0\bigr)
\]

这个损失本质上是一个带间隔的度量学习 / 排序损失，可以看成是 triplet loss 的一种变体。

\section{它在比较什么}

先把符号拆开：

\begin{itemize}
  \item \(A(\hat{e}_c)\)：某个样本经过映射器 \(A\) 后的特征
  可以理解为 anchor
  \item \(f_{text}(t_c)\)：正确类别 \(c\) 的文本特征
  可以理解为 positive
  \item \(f_{text}(t_{\cancel{c}})\)：错误类别的文本特征
  可以理解为 negative
  \item \(\text{dist}(\cdot,\cdot)\)：距离函数
  \item \(m\)：margin，间隔
\end{itemize}

所以它在做的事是：

让 anchor 到正确类别文本的距离，比到错误类别文本的距离至少小 m。

也就是希望满足：

\[
\text{dist}(A(\hat{e}_c), f_{text}(t_c)) + m
\;<\;
\text{dist}(A(\hat{e}_c), f_{text}(t_{\cancel{c}}))
\]

\section{条件已经满足时}

如果已经有：

\[
\text{dist}(A(\hat{e}_c), f_{text}(t_c)) - \text{dist}(A(\hat{e}_c), f_{text}(t_{\cancel{c}})) + m < 0
\]

那么： \( \max(\cdots,0)=0\)

也就是不再惩罚，说明当前样本已经分得足够开了。

\section{条件不满足时}

如果：

\[
\text{dist}(A(\hat{e}_c), f_{text}(t_c)) - \text{dist}(A(\hat{e}_c), f_{text}(t_{\cancel{c}})) + m > 0
\]

说明 anchor 离正确文本还不够近，或者离错误文本还不够远，或者二者之间没有拉开足够的 margin。

这时 loss 就会大于 0，推动模型去更新参数。

\section{它的核心作用}

这个损失把样本特征拉向正确类别的文本特征，同时把它推离错误类别的文本特征，并且要求两者之间有一个安全间隔。

它同时做两件事：一方面让 \(A(\hat{e}_c) \rightarrow f_{text}(t_c)\)，另一方面让 \(A(\hat{e}_c) \text{ 远离 } f_{text}(t_{\cancel{c}})\)。

\section{为什么要加 margin \(m\)}

如果没有 margin，只要求：

\[
\text{dist}(A(\hat{e}_c), f_{text}(t_c))
<
\text{dist}(A(\hat{e}_c), f_{text}(t_{\cancel{c}}))
\]

那模型可能只做到“勉强更近一点”，判别边界会很脆弱。

加入 \emph{m} 后，正确类必须明显更近，而不只是勉强更近。margin 让判别边界更稳，也避免不同类别的特征挤在一起。

\section{它和 triplet loss 的关系}

你这个式子其实很像标准 triplet loss：

\[
\mathcal{L}_{triplet}
=
\max\bigl(
d(a,p)-d(a,n)+m,\ 0
\bigr)
\]

对应关系是 \(a = A(\hat{e}_c)\)，\(p = f_{text}(t_c)\)，\(n = f_{text}(t_{\cancel{c}})\)。

所以它可以理解成一个以“样本特征”为 anchor、以“正确/错误类别文本特征”为正负样本的 triplet-style loss。

\section{在训练中它具体会带来什么效果}

如果模型优化成功，属于类别 \(c\) 的样本会靠近文本原型 \(f_{text}(t_c)\)，远离其他类别文本原型，类别之间的边界也更清晰。

所以它适合类别语义对齐、图文共同空间学习、原型约束，以及零样本或增量学习中的语义锚定。

\section{这个损失相比 CE loss 的特点}

和普通交叉熵相比，它更强调相对几何关系，而不是单纯分类概率。交叉熵关心“哪个 logit 最大”，hinge loss 关心“正确类和错误类之间的距离差有没有超过 margin”。

所以它适合学习结构化特征空间、做类间分离，以及和文本原型或类原型结合。
